% ============================================================
% Section 1: Introduction & Motivation
% ============================================================

\section{Introduction}
\label{sec:intro}

Voice interfaces in automotive environments have become increasingly prevalent,
with modern vehicles integrating speech-controlled infotainment systems,
navigation aids, and climate controls. Current commercial assistants—
including Apple CarPlay's Siri, Google Assistant for Android Auto, and Amazon
Alexa Auto—are designed for a \emph{single active speaker} issuing commands
at a time. In practice, however, a car cabin hosts multiple occupants who may
speak simultaneously, issue overlapping commands, or engage in natural
conversation while incidentally directing requests at the vehicle's assistant.
No public benchmark or open-source system currently addresses this
\emph{in-car multi-speaker ASR} scenario comprehensively.

This paper presents the first end-to-end pipeline for \textbf{In-Car
Multi-Speaker Automatic Speech Recognition} (InCar-MSAR), evaluated on the
newly released AISHELL-5 dataset~\cite{aishell5_2025}—the first large-scale,
publicly available corpus of Mandarin Chinese multi-speaker conversations
recorded in real automotive environments using a four-microphone array. Our
system receives a four-channel audio stream (16\,kHz, one microphone per
door) and produces, for each detected speaker: (i) a clean speech
transcript, (ii) a speaker role label (Driver or Passenger), and (iii) a
structured intent representation suitable for vehicle control.

\subsection*{Contributions}
The main contributions of this work are:
\begin{enumerate}
    \item \textbf{First public benchmark on AISHELL-5}: We report
    Word Error Rate (WER), Concatenated Permutation WER (cpWER), and
    Speaker Attribution Accuracy on both evaluation sets, establishing
    strong baselines for future research.

    \item \textbf{Multi-channel separation pipeline}: We demonstrate that
    combining neural beamforming via SepFormer~\cite{subakan2021attention}
    with large-scale pre-trained ASR (Whisper~\cite{radford2023robust})
    yields statistically significant WER reduction compared to
    single-channel direct transcription.

    \item \textbf{Reproducible open-source system}: The complete pipeline—
    data loading, separation, ASR, speaker role classification, and intent
    understanding—is released as a reproducible Python package with a
    Streamlit demo, Docker image, and evaluation scripts.
\end{enumerate}

\noindent The remainder of this paper is structured as follows:
Section~\ref{sec:related} reviews related work;
Section~\ref{sec:methodology} describes the dataset and proposed architecture;
Section~\ref{sec:experiments} presents experimental results and ablation
studies; Section~\ref{sec:conclusion} concludes with future directions.
